\documentclass{article}
\usepackage[utf8]{inputenc}
\usepackage[brazil]{babel}
\usepackage{geometry}
\geometry{a4paper, left=2cm, right=2cm, top=2cm, bottom=2cm}
\usepackage{hyperref}
\hypersetup{
    colorlinks=true,
    linkcolor=blue,
    filecolor=magenta,      
    urlcolor=cyan,
}

\title{Aula: Identidade Digital}
\author{Professor(a): Seu Nome}
\date{}

\begin{document}

\maketitle

\section*{Objetivos da Aula}
Ao final da aula, os alunos deverão:
\begin{itemize}
    \item Compreender o conceito de identidade digital.
    \item Reconhecer a importância de gerenciar sua presença online.
    \item Identificar riscos e oportunidades relacionados à identidade digital.
    \item Aplicar boas práticas para construir uma identidade digital positiva.
\end{itemize}

\section*{1. Introdução (\textbf{10 minutos})}
\subsection*{O que é Identidade Digital?}
\begin{itemize}
    \item \textbf{Definição:} Conjunto de informações sobre uma pessoa disponíveis na internet (redes sociais, perfis profissionais, comentários, fotos, etc.).
    \item \textbf{Comparação:} Assim como temos uma identidade no mundo físico (RG, CPF), temos uma no mundo virtual.
    \item \textbf{Exemplos:} Perfil no LinkedIn, postagens no Instagram, histórico de buscas no Google.
\end{itemize}

\subsection*{Atividade Inicial (Discussão)}
\textbf{Pergunta aos alunos:} 
\begin{quote}
    ``O que você acha que sua identidade digital diz sobre você?''
\end{quote}

\section*{2. Desenvolvimento (\textbf{30 minutos})}
\subsection*{2.1. Componentes da Identidade Digital}
\begin{itemize}
    \item \textbf{Perfis sociais} (Facebook, Instagram, TikTok)
    \item \textbf{Presença profissional} (LinkedIn, Lattes)
    \item \textbf{Registros públicos} (comentários em fóruns, avaliações)
    \item \textbf{Dados implícitos} (histórico de navegação, cookies)
\end{itemize}

\subsection*{2.2. Riscos e Oportunidades}
\begin{itemize}
    \item \textbf{Riscos:}
        \begin{itemize}
            \item Exposição excessiva (\textit{oversharing})
            \item Cyberbullying e assédio online
            \item Roubo de dados e fraudes
        \end{itemize}
    \item \textbf{Oportunidades:}
        \begin{itemize}
            \item Networking profissional
            \item Construção de marca pessoal
            \item Acesso a oportunidades educacionais e de emprego
        \end{itemize}
\end{itemize}

\subsection*{2.3. Boas Práticas}
\begin{itemize}
    \item \textbf{Privacidade:} Configurar permissões em redes sociais.
    \item \textbf{Reputação:} Pensar antes de postar (``Se não diria pessoalmente, não poste'').
    \item \textbf{Segurança:} Usar senhas fortes e autenticação em dois fatores.
    \item \textbf{Curadoria:} Manter perfis profissionais atualizados (LinkedIn).
\end{itemize}

\section*{3. Atividade Prática (\textbf{10 minutos})}
\subsection*{Cenário para Debate}
Dividir a turma em grupos e discutir:
\begin{quote}
    ``Um recrutador verificou as redes sociais de um candidato e encontrou fotos inadequadas. O candidato foi reprovado. Isso é justo?''
\end{quote}

\section*{4. Conclusão (\textbf{5 minutos})}
\subsection*{Recapitulação}
\begin{itemize}
    \item Identidade digital é o seu ``eu virtual'' e impacta sua vida real.
    \item É possível gerenciá-la com consciência e segurança.
\end{itemize}

\subsection*{Tarefa para Casa}
\begin{itemize}
    \item Pesquisar sobre \href{https://digitalguardian.com/blog/what-digital-footprint}{pegada digital} e trazer exemplos na próxima aula.
\end{itemize}

\end{document}
